Vision-language-action models (VLAs) fine-tune large language models on robot demonstrations, but where the representations driving action prediction originate remains unknown.
We perform the first representational decomposition of a VLA (OpenVLA-7B on LIBERO), separately quantifying action-relevant structure at the vision encoder, projector, and each LLM layer using Representational Similarity Analysis and functional probes.
Three findings emerge.
First, self-supervised pretraining provides the strongest action-relevant geometry: standalone DINOv2 achieves RSA $\rho = 0.692$, exceeding all VL-pretrained encoders (DinoSigLIP+Projector $\rho = 0.615$; CLIP $\rho = 0.557$) and the VLA's own internal representations.
Second, the LLM applies a non-monotonic transformation---the first 12 layers \emph{enhance} geometric alignment with the action space above the vision encoder input ($\rho$: $0.612 \to 0.628$ at L12), while the remaining 20 layers degrade it ($0.628 \to 0.422$)---and concentrates a functionally superior representation at the readout token (probe $R^2 = 0.678$ vs.\ $0.630$ for the projector).
Third, training produces at most a modest geometric gain (RSA $\Delta = 0.017$--$0.023$ across two benchmarks, $3$--$4\%$ of baseline) while substantially improving functional capacity ($R^2$ $\Delta = +0.015$ to $+0.069$, $p \leq 0.006$)---redistribution dominates over creation.
These results indicate that action representations in VLAs are inherited from pretrained vision encoders and non-monotonically transformed by the LLM, which routes instruction identity through the readout token (100\% classification from L1) while preserving visual geometry in image tokens.
Vision encoder quality is the primary design lever; the LLM functions as a routing and preservation mechanism rather than a source of new action-relevant structure.
